\chapter{PSY Ledger --- Psychology Tags}
\label{chap:PSY_Ledger}

\section{Purpose}

The PSY ledger captures psychological structures, affective states,
motivation patterns, and identity-relevant cognitive dynamics.

\section{Scope}

Tags describe:
\begin{itemize}
    \item emotion and affect,
    \item attachment styles,
    \item drives and needs,
    \item defensive patterns,
    \item archetypal behaviors,
    \item pathological conditions (CME-safe framing).
\end{itemize}

\section{Schema}

\[
T_{\text{PSY}} = \langle key,\, gloss,\, bindings,\, constraints \rangle
\]

Examples:
\begin{itemize}
    \item \texttt{psy:arousal}
    \item \texttt{psy:defense-mechanism}
    \item \texttt{psy:projection}
    \item \texttt{psy:archetype-shadow}
    \item \texttt{psy:integration}
\end{itemize}

\section{CME/Jung/Freud Note}

\begin{itemize}
    \item Jungian tags refer to archetypes, shadow, anima/animus, individuation.
    \item Freudian tags refer to ego defenses, drives, sublimation, conflict.
    \item CME-pathology tags should avoid real clinical labels unless marked
          as metaphoric or narrative-contextual.
\end{itemize}

\section{Class Binding Notes}

\begin{itemize}
    \item L-class and O-class dominant.
    \item C-class for stable psychological commitments only.
    \item P-class requires HITL review (identity-impacting).
\end{itemize}

\section{Constraints}

\begin{itemize}
    \item May not assign mental illness labels without operator mediation.
    \item Archetypal tags must be descriptive, not deterministic.
    \item Identity-impacting psychological assertions require HITL.
\end{itemize}
